% math0 manuscript. Build with: latexmk -pdf paper.tex
\documentclass[11pt]{amsart}
\usepackage[margin=1.1in]{geometry}
\usepackage{amsmath,amssymb,amsthm,mathtools}
\usepackage{url}
\usepackage{hyperref}
\usepackage{xcolor}
\usepackage{listings}
\usepackage{newunicodechar}

\lstdefinelanguage{lean}{
  keywords={theorem,lemma,def,abbrev,example,fun,by,have,show,from,let,in,if,then,else,match,with,
    namespace,end,open,import,section,variable,instance,structure,where,calc,obtain,intro,exact,
    refine,rcases,rintro,simp,omega,decide,norm_num,linarith,nlinarith,positivity,induction,cases,
    private,rw,linear_combination},
  sensitive=true, comment=[l]{--}, morecomment=[s]{/-}{-/}, string=[b]"}
\lstset{language=lean, basicstyle=\ttfamily\footnotesize, columns=fullflexible, keepspaces=true, breaklines=true,
  keywordstyle=\color{blue!60!black}\bfseries, commentstyle=\color{gray}, frame=single,
  extendedchars=true, inputencoding=utf8,
  literate={ℕ}{{$\mathbb{N}$}}1 {ℤ}{{$\mathbb{Z}$}}1 {ℚ}{{$\mathbb{Q}$}}1 {∀}{{$\forall$}}1
    {∃}{{$\exists$}}1 {↔}{{$\leftrightarrow$}}1 {¬}{{$\neg$}}1 {∧}{{$\wedge$}}1 {∨}{{$\vee$}}1
    {→}{{$\to$}}1 {≤}{{$\le$}}1}
% Unicode commonly found in Lean sources.
\newunicodechar{ℕ}{\ensuremath{\mathbb{N}}}
\newunicodechar{ℤ}{\ensuremath{\mathbb{Z}}}
\newunicodechar{ℚ}{\ensuremath{\mathbb{Q}}}
\newunicodechar{ℝ}{\ensuremath{\mathbb{R}}}
\newunicodechar{ℂ}{\ensuremath{\mathbb{C}}}
\newunicodechar{∀}{\ensuremath{\forall}}
\newunicodechar{∃}{\ensuremath{\exists}}
\newunicodechar{→}{\ensuremath{\to}}
\newunicodechar{←}{\ensuremath{\leftarrow}}
\newunicodechar{↔}{\ensuremath{\leftrightarrow}}
\newunicodechar{↦}{\ensuremath{\mapsto}}
\newunicodechar{¬}{\ensuremath{\neg}}
\newunicodechar{∧}{\ensuremath{\wedge}}
\newunicodechar{∨}{\ensuremath{\vee}}
\newunicodechar{≤}{\ensuremath{\le}}
\newunicodechar{≥}{\ensuremath{\ge}}
\newunicodechar{≠}{\ensuremath{\ne}}
\newunicodechar{∈}{\ensuremath{\in}}
\newunicodechar{∉}{\ensuremath{\notin}}
\newunicodechar{⊆}{\ensuremath{\subseteq}}
\newunicodechar{∩}{\ensuremath{\cap}}
\newunicodechar{∪}{\ensuremath{\cup}}
\newunicodechar{∑}{\ensuremath{\sum}}
\newunicodechar{∏}{\ensuremath{\prod}}
\newunicodechar{⟨}{\ensuremath{\langle}}
\newunicodechar{⟩}{\ensuremath{\rangle}}
\newunicodechar{·}{\ensuremath{\cdot}}
\newunicodechar{∘}{\ensuremath{\circ}}
\newunicodechar{⁻}{\ensuremath{^{-}}}
\newunicodechar{¹}{\ensuremath{^{1}}}
\newunicodechar{²}{\ensuremath{^{2}}}
\newunicodechar{₀}{\ensuremath{_{0}}}
\newunicodechar{₁}{\ensuremath{_{1}}}
\newunicodechar{₂}{\ensuremath{_{2}}}
\newunicodechar{α}{\ensuremath{\alpha}}
\newunicodechar{β}{\ensuremath{\beta}}
\newunicodechar{ε}{\ensuremath{\varepsilon}}
\newunicodechar{δ}{\ensuremath{\delta}}
\newunicodechar{π}{\ensuremath{\pi}}
\newunicodechar{σ}{\ensuremath{\sigma}}
\newunicodechar{φ}{\ensuremath{\varphi}}
\newunicodechar{∞}{\ensuremath{\infty}}
\newunicodechar{⌊}{\ensuremath{\lfloor}}
\newunicodechar{⌋}{\ensuremath{\rfloor}}
\newunicodechar{⌈}{\ensuremath{\lceil}}
\newunicodechar{⌉}{\ensuremath{\rceil}}
\newunicodechar{∣}{\ensuremath{\mid}}
\newunicodechar{≡}{\ensuremath{\equiv}}
\newunicodechar{×}{\ensuremath{\times}}
\newunicodechar{⧸}{\ensuremath{/}}
\newunicodechar{‖}{\ensuremath{\|}}
\newunicodechar{√}{\ensuremath{\surd}}
\newunicodechar{∅}{\ensuremath{\emptyset}}
\newunicodechar{⦃}{\ensuremath{\{\!|}}
\newunicodechar{⦄}{\ensuremath{|\!\}}}
\newunicodechar{𝓝}{\ensuremath{\mathcal{N}}}
\newunicodechar{≫}{\ensuremath{\gg}}
\newunicodechar{≪}{\ensuremath{\ll}}

\theoremstyle{plain}
\newtheorem{theorem}{Theorem}
\newtheorem{lemma}[theorem]{Lemma}
\newtheorem{proposition}[theorem]{Proposition}
\newtheorem{corollary}[theorem]{Corollary}
\theoremstyle{definition}
\newtheorem{definition}[theorem]{Definition}
\newtheorem{remark}[theorem]{Remark}
\newtheorem{example}[theorem]{Example}

\emergencystretch=3em
% Lean identifiers in running text (breakable at dots and underscores).
\expandafter\def\expandafter\UrlBreaks\expandafter{\UrlBreaks\do\_\do\.}
\DeclareUrlCommand\decl{\urlstyle{tt}}
\newcommand{\N}{\mathbb{N}}
\newcommand{\Z}{\mathbb{Z}}
\newcommand{\Q}{\mathbb{Q}}
\newcommand{\Tot}{\operatorname{Tot}}
\newcommand{\R}{\mathcal{R}}
\newcommand{\cat}{\mathbin{+\!\!+}}
\newcommand{\iv}[2]{[#1\,..\,#2]}

\title[Neder's conjecture for OEIS A308121]{For odd $n>1$, row $2n$ of the triangle $T(n,k)=\operatorname{den}\bigl(\tfrac{\varphi(n)}{n}\bigr)k-\operatorname{num}\bigl(\tfrac{\varphi(n)}{n}\bigr)r_k(n)$ is row $n$ with its halves interchanged: Neder's conjecture for OEIS A308121}
\author{math0}
\address{Produced by Claude (an AI model developed by Anthropic) in the math0 project}
\date{October 25, 2026}

\begin{document}
\begin{abstract}
For $n\ge1$ let $r_1(n)<\dots<r_{\varphi(n)}(n)$ be the integers in $[1,n]$ coprime to $n$, and write $\varphi(n)/n=a_n/b_n$ in lowest terms; OEIS A308121 is the irregular triangle whose $n$-th row is $\bigl(b_nk-a_nr_k(n)\bigr)_{1\le k\le\varphi(n)}$.
We prove the conjecture recorded in this entry by Charlie Neder in 2019 that for odd $n>1$, row $2n$ is obtained from row $n$ by interchanging its first and second halves.
The proof lists the totatives of $2n$ as the images of the totatives of $n$ above $n/2$ under $r\mapsto 2r-n$, followed by the images of those below $n/2$ under $r\mapsto n+2r$, and shows that $\varphi(2n)/(2n)$ has reduced form $(a_n/2)/b_n$.
To our knowledge the result is new, although the argument is elementary; the OEIS entry still lists the statement as a conjecture.
The theorem is verified in Lean~4 against a frozen, fidelity-reviewed statement that math0 formalized from the OEIS entry, whose definitions are checked by the kernel against the $72$ listed terms; this statement is not part of the formal-conjectures repository.
\end{abstract}
\maketitle

\section{Introduction}

For an integer $n\ge1$ let $\varphi$ be Euler's totient function and let
\[
  r_1(n)<r_2(n)<\dots<r_{\varphi(n)}(n)
\]
be the integers $r$ with $1\le r\le n$ and $\gcd(r,n)=1$; this is row $n$ of the triangle A038566 \cite{A038566}, so that row $1$ is $(1)$. Write
\[
  \frac{\varphi(n)}{n}=\frac{a_n}{b_n},\qquad \gcd(a_n,b_n)=1,\quad b_n\ge1,
\]
so that $b_n$ is A109395$(n)$ \cite{A109395} and $a_n$ is A076512$(n)$ \cite{A076512}.
The sequence A308121 \cite{A308121}, contributed by Jamie Morken on May~13, 2019, is the irregular triangle
\[
  T(n,k)=b_n\,k-a_n\,r_k(n)\qquad(n\ge1,\ 1\le k\le\varphi(n)),
\]
whose rows $1,\dots,15$ are
\[
\begin{array}{r|l@{\qquad}r|l@{\qquad}r|l}
 1&0 & 6&2,1 & 11&1,2,\dots,10\\
 2&1 & 7&1,2,3,4,5,6 & 12&2,1,2,1\\
 3&1,2 & 8&1,1,1,1 & 13&1,2,\dots,12\\
 4&1,1 & 9&1,2,1,2,1,2 & 14&4,5,6,1,2,3\\
 5&1,2,3,4 & 10&3,4,1,2 & 15&7,14,13,4,11,2,1,8
\end{array}
\]
The entry records several properties of the rows; among them, Charlie Neder added on June~5, 2019 the observation that if a prime $p$ divides $n$ then row $pn$ consists of $p$ copies of row $n$, together with the following conjecture: \emph{if $n$ is odd, then row $2n$ can be obtained from row $n$ by interchanging the first and second halves.}
The table illustrates it for $n=3,5,7$ ($6,10,14$); for $n=1$ it cannot hold, since row $1$ has a single entry.
We prove the conjecture for every odd $n>1$ (Theorem~\ref{thm:main}).
We do not treat Neder's observation on rows $pn$.

\subsection*{Method}
For odd $n\ge3$ the totatives of $2n$ are, in increasing order, $2r-n$ for the totatives $r>n/2$ of $n$, followed by $n+2r$ for the totatives $r<n/2$, and exactly half of the totatives of $n$ lie on each side of $n/2$ (Lemmas~\ref{lem:tot} and~\ref{lem:tot2}). Moreover $\varphi(2n)/(2n)$ has reduced form $(a_n/2)/b_n$ (Lemma~\ref{lem:frac}). Substituting both facts into the definition of $T(2n,k)$ gives the claim.

\subsection*{Status and prior work}
The statement is still marked ``Conjecture'' in A308121 (revision \#47, August~5, 2019), with no later comment or edit note recording a proof or a counterexample.
Two other OEIS entries cite A308121: A307964 \cite{A307964} (J.~Morken, July~2019), which restricts the triangle to the rows of odd squarefree composite numbers and calls them its ``primitive rows'', and A309497 \cite{A309497} (J.~Morken, August~2019), which collects the ``Primorial rows of A308121''. Neither addresses the conjecture, and the entries A038566, A076512 and A109395 used in the definition say nothing about it.
The math0 referees found no proof or counterexample on arXiv, MathOverflow, Mathematics Stack Exchange or the OEIS wiki; the SeqFan mailing-list archive could not be reached.
The conjecture is not among the OEIS statements of the formal-conjectures repository \cite{FC}, and hence not among the $492$ conjectures of the OEIS Open benchmark \cite{Ad}, whose Lean statements (formalized by Tsoukalas et al.) were released in that repository.
The ingredients are classical: $\varphi(2n)=\varphi(n)$ for odd $n$, and the totatives of $n$ and of $2n$ correspond bijectively. For instance, Nijenhuis \cite{Nij} uses the group isomorphism $y\mapsto y$ ($y$ odd), $y\mapsto y+n$ ($y$ even) from the reduced residues modulo $n$ to those modulo $2n$, whose proof he calls a simple exercise; our bijection $r\mapsto(2r+n)\bmod 2n$ is this isomorphism composed with doubling modulo $n$ (Remark~\ref{rem:map}).
We therefore classify the result as \emph{new}, with the caveat that the argument is elementary and may well have been noticed without being recorded.

\section{Statement}\label{sec:statement}

\subsection*{Notation}
All lists are finite sequences of integers. For integers $x\le y+1$ we write $\iv{x}{y}$ for the increasing list $(x,x+1,\dots,y)$ (empty if $x=y+1$); $\ell\cat\ell'$ is concatenation, $|\ell|$ the length, $\ell_i$ the $i$-th entry ($1\le i\le|\ell|$), $f(\ell)$ the list obtained by applying $f$ to every entry, and $\ell|_P$ the sublist of the entries satisfying a predicate $P$, in their original order.
For $n\ge1$ let $C_n(r)$ be the predicate $\gcd(r,n)=1$ and
\[
  \Tot(n)=\iv{1}{n}|_{C_n}=(r_1(n),\dots,r_{\varphi(n)}(n)),\qquad
  \R(n)=\bigl(T(n,1),\dots,T(n,\varphi(n))\bigr),
\]
so that $\R(n)$ is row $n$ of A308121.

The COMMENTS section of A308121 \cite{A308121} contains, verbatim:
\begin{quote}
{\small\texttt{From \_Charlie Neder\_, Jun 05 2019: (Start)\\
If p is a prime dividing n, then row p*n consists of p copies of row n.\\
Conjecture: If n is odd, then row 2n can be obtained from row n by interchanging the first and second halves. (End)}}
\end{quote}

\begin{theorem}\label{thm:main}
Let $n>1$ be odd. Then $\varphi(n)=\varphi(2n)=2h$ for an integer $h\ge1$, and
\[
  T(2n,k)=T(n,h+k),\qquad T(2n,h+k)=T(n,k)\qquad(1\le k\le h).
\]
Equivalently, there are lists $A,B$ with $|A|=|B|$, $\R(n)=A\cat B$ and $\R(2n)=B\cat A$.
\end{theorem}

For $n\ge3$ the number $\varphi(n)$ is even, so the halves of $\R(n)$ exist and are unique; $\R(1)=(0)$ has no halves.
For even $n$ the statement does not apply, since then $\varphi(2n)=2\varphi(n)$.

\subsection*{Lean formalization}
The statement was formalized by math0 in the file \path{lean/Math0/Statements/OEIS/A308121.lean}. It was reviewed for fidelity and frozen (SHA-256 \texttt{70da202a\ldots c9e58}) before any proof attempt. Its definitions read as follows (verbatim; the copyright header, the import, the module docstring, the docstrings, blank lines and the four theorems are omitted):
\begin{lstlisting}
namespace Math0.Statements.A308121
def totatives (n : ℕ) : List ℕ :=
  (List.range' 1 n).filter (fun r => Nat.Coprime r n)
def A109395 (n : ℕ) : ℕ := ((Nat.totient n : ℚ) / n).den
def A076512 (n : ℕ) : ℤ := ((Nat.totient n : ℚ) / n).num
def T (n k : ℕ) : ℤ :=
  (A109395 n : ℤ) * k - A076512 n * ((totatives n).getD (k - 1) 0 : ℕ)
def row (n : ℕ) : List ℤ :=
  (List.range (Nat.totient n)).map (fun i => T n (i + 1))
def statement : Prop :=
  ∀ n : ℕ, Odd n → 1 < n →
    ∃ A B : List ℤ, row n = A ++ B ∧ A.length = B.length ∧ row (2 * n) = B ++ A
end Math0.Statements.A308121
\end{lstlisting}
Here \texttt{List.range' 1 n} is $\iv{1}{n}$, so \decl{totatives} is $\Tot(n)$, i.e.\ row $n$ of A038566, described there as the ``triangle in which n-th row gives phi(n) numbers between 1 and n that are relatively prime to n'' \cite{A038566}; in particular \texttt{totatives 1 = [1]}.
\decl{A109395} and \decl{A076512} are the denominator and numerator of the reduced rational number $\varphi(n)/n$, as in the names of A109395 and (via its comment ``Numerator of phi(n)/n'') A076512; Lean's \texttt{Rat.num} and \texttt{Rat.den} return the reduced form with positive denominator.
\decl{T} computes $b_nk-a_nr_k(n)$ in $\Z$, because the sequence is signed; the index $k$ is $1$-based through \texttt{getD (k - 1) 0} (only the values with $1\le k\le\varphi(n)$ are used), and \decl{row} lists $T(n,1),\dots,T(n,\varphi(n))$, matching the comment ``Row n has length A000010(n)''.
The conjecture is encoded literally: the first and second halves of row $n$ are lists $A,B$ of equal length with $\R(n)=A\cat B$, and $\R(2n)=B\cat A$.
The hypothesis \texttt{1 < n} excludes $n=1$, for which row $1=(0)$ has no halves (and row $2=(1)$); the statement file documents this reading.
Four theorems in the statement file, all proved by \texttt{decide +kernel} (evaluation in the kernel), compare the definitions with the OEIS entry; two of them read:
\begin{lstlisting}
theorem a_data : (List.range 15).flatMap (fun i => row (i + 1)) =
    [0, 1, 1, 2, 1, 1, 1, 2, 3, 4, 2, 1, 1, 2, 3, 4, 5, 6, 1, 1, 1, 1, 1, 2, 1, 2, 1, 2, 3, 4,
      1, 2, 1, 2, 3, 4, 5, 6, 7, 8, 9, 10, 2, 1, 2, 1, 1, 2, 3, 4, 5, 6, 7, 8, 9, 10, 11, 12,
      4, 5, 6, 1, 2, 3, 7, 14, 13, 4, 11, 2, 1, 8]
theorem T_210_2 : T 210 2 = -18
\end{lstlisting}
\decl{a_data} checks that rows $1,\dots,15$, flattened, are exactly the $72$ listed terms $a(1),\dots,a(72)$ of A308121; \decl{T_210_2} checks the worked example $T(210,2)=35\cdot2-8\cdot11=-18$ of the entry; \decl{statement_check_7} and \decl{statement_check_15} check the halves swap for $n=7$ and $n=15$ (sanity checks of the transcription, not part of the proof).

\section{Proof}\label{sec:proof}

The proof follows the Lean file \path{lean/Math0/Frontier/OEIS/A308121.lean}; the names of the corresponding Lean declarations are given with each step.
Throughout, $t\ge1$ and $n=2t+1$, and
\[
  L=\iv{1}{t}|_{C_n},\qquad H=\iv{t+1}{2t}|_{C_n},\qquad h=|L|.
\]
Thus $L$ lists the totatives of $n$ below $n/2$ and $H$ those above $n/2$, both increasing.

\begin{lemma}[\decl{len_tot}]\label{lem:len}
For every $m\ge1$, $|\Tot(m)|=\varphi(m)$.
\end{lemma}

\begin{proof}
$|\Tot(m)|$ is the number of integers $r\in[1,m]$ with $\gcd(r,m)=1$, which is $\varphi(m)$. In Lean this is Mathlib's \decl{Nat.filter_coprime_Ico_eq_totient}, which counts the integers coprime to $m$ in any interval of $m$ consecutive integers.
\end{proof}

\begin{lemma}[\decl{tot_split}, \decl{map_sub_range'}, \decl{len_eq}]\label{lem:tot}
$\Tot(n)=L\cat H$ and $|H|=|L|=h$. Consequently $\varphi(n)=2h$.
\end{lemma}

\begin{proof}
We have $\iv{1}{n}=\iv{1}{t}\cat\iv{t+1}{2t}\cat(n)$, and $\gcd(n,n)=n>1$, so $n$ is removed by the filter. Filtering commutes with concatenation, so $\Tot(n)=L\cat H$.

Let $\rho(r)=n-r$. Then $\rho\bigl(\iv{1}{t}\bigr)=(n-1,n-2,\dots,n-t)$ is the reverse of $\iv{t+1}{2t}$ (\decl{map_sub_range'}). Since $\gcd(n-r,n)=\gcd(r,n)$ for $1\le r\le t$, the predicates $C_n\circ\rho$ and $C_n$ agree on $\iv{1}{t}$. Filtering commutes with reversal, and the filter of a mapped list is the map of the filter by the composed predicate, so
\[
  \operatorname{rev}(H)=\iv{t+1}{2t}^{\mathrm{rev}}\big|_{C_n}=\rho\bigl(\iv{1}{t}\bigr)\big|_{C_n}=\rho\bigl(\iv{1}{t}|_{C_n\circ\rho}\bigr)=\rho(L).
\]
Hence $|H|=|L|=h$, and $\varphi(n)=|L|+|H|=2h$ by Lemma~\ref{lem:len}.
\end{proof}

\begin{lemma}[\decl{filter_pair}]\label{lem:pair}
Let $m\ge1$, let $a$ be odd and $k\ge0$. Then
\[
  \iv{a}{a+2k-1}\big|_{C_{2m}}=(a,a+2,\dots,a+2k-2)\big|_{C_m}.
\]
\end{lemma}

\begin{proof}
Induction on $k$; the case $k=0$ is trivial. The list $\iv{a}{a+2k+1}$ begins with $a,a+1$, followed by $\iv{a+2}{a+2k+1}$, and $a+2$ is odd. Since $a$ is odd, $\gcd(a,2m)=1$ if and only if $\gcd(a,m)=1$; since $a+1$ is even, $\gcd(a+1,2m)\ne1$. Applying the induction hypothesis to $a+2$ gives the claim for $k+1$.
\end{proof}

\begin{lemma}[\decl{range'_two_eq_map}, \decl{cop_hi}, \decl{cop_lo}, \decl{tot2_split}]\label{lem:tot2}
Let $\sigma(r)=2r-n$ and $\tau(r)=n+2r$. Then
\[
  \Tot(2n)=\sigma(H)\cat\tau(L).
\]
Consequently $\varphi(2n)=2h$.
\end{lemma}

\begin{proof}
Split $\iv{1}{2n}=\iv{1}{2t}\cat(2t+1,2t+2)\cat\iv{2t+3}{4t+2}$. The middle entries are removed by the filter $C_{2n}$: $\gcd(2t+1,2n)=n>1$, and $2t+2$ is even. By Lemma~\ref{lem:pair} with $m=n$, $k=t$ and $a=1$, respectively $a=2t+3$,
\[
  \Tot(2n)=(1,3,\dots,2t-1)\big|_{C_n}\cat(2t+3,2t+5,\dots,4t+1)\big|_{C_n}.
\]
Since $\sigma(t+1)=1$ and $\sigma(r+1)=\sigma(r)+2$, we have $(1,3,\dots,2t-1)=\sigma\bigl(\iv{t+1}{2t}\bigr)$; likewise $\tau(1)=2t+3$ and $\tau(r+1)=\tau(r)+2$ give $(2t+3,\dots,4t+1)=\tau\bigl(\iv{1}{t}\bigr)$ (\decl{range'_two_eq_map}).
Because $n$ is odd, $\gcd(2r-n,n)=\gcd(2r,n)=\gcd(r,n)$ for $2r\ge n$ (\decl{cop_hi}) and $\gcd(n+2r,n)=\gcd(2r,n)=\gcd(r,n)$ (\decl{cop_lo}). Hence $C_n\circ\sigma=C_n$ on $\iv{t+1}{2t}$ and $C_n\circ\tau=C_n$ on $\iv{1}{t}$, and as in the proof of Lemma~\ref{lem:tot},
\[
  \sigma\bigl(\iv{t+1}{2t}\bigr)\big|_{C_n}=\sigma(H),\qquad \tau\bigl(\iv{1}{t}\bigr)\big|_{C_n}=\tau(L).
\]
Finally $\varphi(2n)=|\sigma(H)|+|\tau(L)|=|H|+|L|=2h$ by Lemmas~\ref{lem:len} and~\ref{lem:tot}.
\end{proof}

In Lean, $\sigma$ is written with truncated subtraction \texttt{2 * r - (2 * t + 1)}; on $\iv{t+1}{2t}$ we have $2r\ge2t+2>n$, so the subtraction is exact.

\begin{lemma}[\decl{rat_facts}]\label{lem:frac}
Let $n>1$ be odd. Then $b_{2n}=b_n$, $2a_{2n}=a_n$, and $a_n\,n=b_n\,\varphi(n)$.
\end{lemma}

\begin{proof}
Let $q=\varphi(n)/n=a_n/b_n$. Clearing denominators gives $a_n\,n=b_n\,\varphi(n)$. Since $\gcd(2,n)=1$, $\varphi(2n)=\varphi(2)\varphi(n)=\varphi(n)$, so $\varphi(2n)/(2n)=q/2$.
As $n>2$, $\varphi(n)$ is even, so $a_n\,n$ is even; as $n$ is odd, $a_n$ is even, say $a_n=2c$. From $\gcd(a_n,b_n)=1$ we get $\gcd(c,b_n)=1$. Therefore $q/2=c/b_n$ with $b_n\ge1$ and $\gcd(c,b_n)=1$, and by uniqueness of the reduced form (Mathlib's \decl{Rat.num_div_eq_of_coprime} and \decl{Rat.den_div_eq_of_coprime}) $a_{2n}=c$ and $b_{2n}=b_n$.
\end{proof}

\begin{proof}[Proof of Theorem~\ref{thm:main} (\decl{result})]
Let $n=2t+1$ with $t\ge1$, and let $L,H,h$ be as above. By Lemmas~\ref{lem:tot} and~\ref{lem:tot2}, $\varphi(n)=\varphi(2n)=2h$. Write $a=a_n$, $b=b_n$ and $c=a_{2n}$. By Lemma~\ref{lem:frac}, $b_{2n}=b$, $a=2c$ and $an=2bh$, so
\begin{equation}\label{eq:key}
  c\,n=b\,h .
\end{equation}
If $\varphi(m)=2h$, then $\R(m)=\bigl(T(m,1),\dots,T(m,h)\bigr)\cat\bigl(T(m,h+1),\dots,T(m,2h)\bigr)$ (\decl{row_split}). Apply this to $m=n$, and let $A$ and $B$ be the two halves of $\R(n)$, so $\R(n)=A\cat B$ and $|A|=|B|=h$. Apply it also to $m=2n$; it remains to show that the first half of $\R(2n)$ is $B$ and the second is $A$, i.e.\ the two identities of the theorem.

Let $1\le k\le h$. By Lemma~\ref{lem:tot2}, $r_k(2n)=\sigma(H_k)=2H_k-n$ and $r_{h+k}(2n)=\tau(L_k)=n+2L_k$; by Lemma~\ref{lem:tot}, $r_k(n)=L_k$ and $r_{h+k}(n)=H_k$. Using $a=2c$ and \eqref{eq:key},
\begin{align*}
  T(2n,k)&=bk-c\,(2H_k-n)=bk-aH_k+cn=b(h+k)-aH_k=T(n,h+k),\\
  T(2n,h+k)&=b(h+k)-c\,(n+2L_k)=b(h+k)-bh-aL_k=bk-aL_k=T(n,k).
\end{align*}
Hence $\R(2n)=B\cat A$.
\end{proof}

In Lean, the index bookkeeping (the positions $k$ and $h+k$ in $L\cat H$ and in $\sigma(H)\cat\tau(L)$) is done with \decl{List.getElem_append_left} and \decl{List.getElem_append_right}, and the two displayed identities are closed by \texttt{linear\_combination} from $2c=a$ and \eqref{eq:key}.

\begin{remark}\label{rem:map}
Lemma~\ref{lem:tot2} says that $r\mapsto(2r+n)\bmod 2n$ maps $\Tot(n)$ bijectively onto $\Tot(2n)$, increasingly on each of $L$ and $H$, with the image of $H$ preceding that of $L$. This map is the composition of doubling modulo $n$, $r\mapsto 2r\bmod n$, with the isomorphism $y\mapsto y$ ($y$ odd), $y\mapsto y+n$ ($y$ even) used by Nijenhuis \cite{Nij}. Theorem~\ref{thm:main} can also be read as $T(2n,k)=T(n,k')$, where $k'\equiv k+h\pmod{2h}$ and $1\le k'\le 2h$.
\end{remark}

\begin{example}
For $n=15$: $\varphi(15)/15=8/15$, $L=(1,2,4,7)$, $H=(8,11,13,14)$, $h=4$, and $\R(15)=(7,14,13,4\mid 11,2,1,8)$. Lemma~\ref{lem:tot2} gives $\Tot(30)=(1,7,11,13,17,19,23,29)$, and $\varphi(30)/30=4/15$, so $c=4$, $b=15$, and $\R(30)=(11,2,1,8\mid 7,14,13,4)$, as in the b-file of A308121 (rows $1,\dots,210$). Similarly row $210$, which contains the entry's worked example $T(210,2)=-18$, is row $105$ with its halves of length $24$ interchanged.
\end{example}

\section{Machine verification}\label{sec:lean}

\subsection*{Files and declarations}
The proof is in the file \path{lean/Math0/Frontier/OEIS/A308121.lean} of the math0 repository, \url{https://github.com/Foundation-0/math0}. It imports the frozen statement file \path{lean/Math0/Statements/OEIS/A308121.lean} (Section~\ref{sec:statement}) and the audit commands in \path{lean/Math0/Audit.lean}. The main declaration is \decl{Math0.Frontier.A308121.result}, of type \decl{Math0.Statements.A308121.statement}. The auxiliary lemmas are private: \decl{filter_pair} (Lemma~\ref{lem:pair}), \decl{range'_two_eq_map}, \decl{map_sub_range'}, \decl{len_tot} (Lemma~\ref{lem:len}), \decl{cop_hi}, \decl{cop_lo}, \decl{tot_split}, \decl{tot2_split}, \decl{len_eq}, \decl{rat_facts} (Lemma~\ref{lem:frac}), \decl{getD_lt} and \decl{row_split}. Two of them, and the end of the file, read verbatim (docstrings and proofs omitted):
\begin{lstlisting}
private lemma tot2_split (t : ℕ) (ht : 1 ≤ t) :
    totatives (2 * (2 * t + 1)) =
      ((List.range' (t + 1) t).filter (fun r => Nat.Coprime r (2 * t + 1))).map
          (fun r => 2 * r - (2 * t + 1)) ++
      ((List.range' 1 t).filter (fun r => Nat.Coprime r (2 * t + 1))).map
          (fun r => (2 * t + 1) + 2 * r)
private lemma rat_facts (n : ℕ) (hn : Odd n) (h1 : 1 < n) :
    A109395 (2 * n) = A109395 n ∧ A076512 (2 * n) * 2 = A076512 n ∧
      A076512 n * n = A109395 n * Nat.totient n
theorem result : statement := by
  ...
end Math0.Frontier.A308121
#math0_settles Math0.Frontier.A308121.result Math0.Statements.A308121.statement
\end{lstlisting}
The proof is general: it contains no finite case check, and it uses no \texttt{decide} or \texttt{native\_decide} tactic and no \texttt{set\_option}. It does not use the data theorems of the statement file.

\subsection*{Audit}
The command \texttt{\#math0\_settles}, defined in \path{lean/Math0/Audit.lean} for statements that math0 formalizes itself, fails unless (1) \decl{statement} is a \texttt{Prop}-valued definition without universe parameters; (2) \decl{statement} and \decl{result} use no axioms besides \decl{propext}, \decl{Classical.choice} and \decl{Quot.sound} (this rejects \decl{sorryAx}, \decl{Lean.ofReduceBool} from \texttt{native\_decide}, and custom axioms); and (3) the type of \decl{result} is syntactically the statement or its negation, up to metadata annotations, with no unfolding. Here it reports
\begin{lstlisting}[basicstyle=\ttfamily\scriptsize]
math0 audit: 'Math0.Statements.A308121.statement' depends on axioms [Classical.choice, Quot.sound, propext]
math0 audit: 'Math0.Frontier.A308121.result' depends on axioms [Classical.choice, Quot.sound, propext]
math0 statement check: 'Math0.Frontier.A308121.result' proves 'Math0.Statements.A308121.statement'
\end{lstlisting}

\subsection*{Toolchain and reproduction}
We use Lean~4 \cite{Lean4} \texttt{v4.33.1} and Mathlib \cite{Mathlib} \texttt{v4.33.1}, pinned in \path{lean/}; the repository also pins formal-conjectures \cite{FC} at commit \texttt{c77fdd4}, which this result does not use.
To reproduce, run \path{scripts/verify.sh} from the repository root: it fetches the Mathlib build cache, builds all math0 results with \texttt{lake build} (re-running every audit command), checks that the pinned \texttt{formal\_conjectures} and \texttt{mathlib} checkouts are unmodified, and reports the number of passed statement checks. The single file can be checked with \texttt{lake env lean Math0/Frontier/OEIS/A308121.lean} in the directory \path{lean}; this takes between about 20 seconds and a minute on our machine, depending on load.
Independently of Lean, the referees checked the conjecture numerically for all odd $n$ with $3\le n<6000$, and the intermediate Lemmas~\ref{lem:tot}, \ref{lem:tot2} and~\ref{lem:frac} for all odd $n$ with $3\le n<4000$.

{\small
\begin{thebibliography}{99}

\bibitem{Ad}
T.~Adamczewski, \emph{OEIS Open: How many conjectures can language models turn into theorems?}, arXiv:2608.11941, 2026.

\bibitem{FC}
M.~Firsching, P.~Lezeau, S.~Mercuri, M.~Z. Horv\'ath, Y.~Dillies, C.~S\"onne, E.~Wieser, F.~Zhang, T.~Hubert, B.~Ag\"uera y Arcas, and P.~Kohli, \emph{Formal Conjectures: An open and evolving benchmark for verified discovery in mathematics}, arXiv:2605.13171, 2026. Repository: \url{https://github.com/google-deepmind/formal-conjectures}.

\bibitem{Mathlib}
The mathlib Community, \emph{The Lean mathematical library}, in: Proceedings of the 9th ACM SIGPLAN International Conference on Certified Programs and Proofs (CPP 2020), ACM, 2020, 367--381.

\bibitem{Lean4}
L.~de Moura and S.~Ullrich, \emph{The Lean 4 theorem prover and programming language}, in: Automated Deduction -- CADE 28, Lecture Notes in Comput. Sci. \textbf{12699}, Springer, 2021, 625--635.

\bibitem{Nij}
A.~Nijenhuis, \emph{Small Gamma products with simple values}, arXiv:0907.1689, 2009.

\bibitem{A038566}
The OEIS Foundation Inc., \url{https://oeis.org/A038566}, entry by N.~J.~A. Sloane.

\bibitem{A076512}
The OEIS Foundation Inc., \url{https://oeis.org/A076512}, entry by R.~Zumkeller (Oct.~15, 2002).

\bibitem{A109395}
The OEIS Foundation Inc., \url{https://oeis.org/A109395}, entry by F.~Vrabec (Aug.~26, 2005).

\bibitem{A307964}
The OEIS Foundation Inc., \url{https://oeis.org/A307964}, entry by J.~Morken (July~29, 2019).

\bibitem{A308121}
The OEIS Foundation Inc., \url{https://oeis.org/A308121}, entry by J.~Morken (May~13, 2019); conjecture by C.~Neder (June~5, 2019); revision \#47, Aug.~5, 2019.

\bibitem{A309497}
The OEIS Foundation Inc., \url{https://oeis.org/A309497}, entry by J.~Morken (Aug.~5, 2019).

\end{thebibliography}
}

\end{document}
